\documentclass[11pt,a4paper]{article}
\usepackage[margin=1in]{geometry}
\usepackage{amsmath,amssymb,amsthm}
\usepackage{algorithm}
\usepackage{algpseudocode}
\usepackage{booktabs}
\usepackage{hyperref}

\theoremstyle{definition}
\newtheorem{definition}{Definition}[section]
\newtheorem{theorem}{Theorem}[section]
\newtheorem{lemma}[theorem]{Lemma}
\newtheorem{remark}{Remark}[section]

\title{\textbf{Rigorous Mathematical Specification, Layer-Wise Composite Mappings, and Universal Approximation of the Multilayer Perceptron (MLP)}}
\author{Team Scaibu \\ \small Email: \texttt{jaydeep.9664920749@gmail.com} \\ \small Architecture and Core Engine Team}
\date{October 2026}

\begin{document}
\maketitle

\begin{abstract}
This document presents the complete theoretical foundation and algorithmic specification of the Multi-Layer Perceptron (MLP) Feedforward Network. By composing alternating affine transformations and non-linear coordinate-wise activation functions across $L$ layers, MLPs form universal function approximators capable of modeling continuous functions on compact domains. We formalize forward layer recurrences, state the Universal Approximation Theorem, formulate complete algorithmic pseudo-code, derive computational complexity, work through a step-by-step numerical example, and discuss stability considerations.
\end{abstract}

\section{Notation and Mathematical Preliminaries}

Let $x \in \mathbb{R}^{d_0}$ denote an input vector, and let $L \in \mathbb{N}_{\ge 1}$ denote the total number of parametric layers in the network.

\begin{itemize}
    \item $d_l \in \mathbb{N}_{\ge 1}$: Dimension of representation space at layer $l \in \{0, 1, \dots, L\}$.
    \item $W^{(l)} \in \mathbb{R}^{d_l \times d_{l-1}}$: Weight matrix for layer $l$.
    \item $b^{(l)} \in \mathbb{R}^{d_l}$: Bias vector for layer $l$.
    \item $z^{(l)} \in \mathbb{R}^{d_l}$: Pre-activation vector at layer $l$.
    \item $a^{(l)} \in \mathbb{R}^{d_l}$: Post-activation vector at layer $l$ ($a^{(0)} \equiv x$).
    \item $\sigma_{\text{hid}}, \sigma_{\text{out}}$: Non-linear activation operators applied at hidden and output layers.
\end{itemize}

\begin{table}[h]
\centering
\caption{Summary of Mathematical Symbols for MLP Feedforward}
\begin{tabular}{lll}
\toprule
\textbf{Symbol} & \textbf{Type / Domain} & \textbf{Description} \\
\midrule
$x$ & $\mathbb{R}^{d_0}$ & Input feature vector \\
$W^{(l)}$ & $\mathbb{R}^{d_l \times d_{l-1}}$ & Layer $l$ weight matrix \\
$b^{(l)}$ & $\mathbb{R}^{d_l}$ & Layer $l$ bias vector \\
$z^{(l)}$ & $\mathbb{R}^{d_l}$ & Pre-activation vector $W^{(l)} a^{(l-1)} + b^{(l)}$ \\
$a^{(l)}$ & $\mathbb{R}^{d_l}$ & Post-activation vector $\sigma(z^{(l)})$ \\
\bottomrule
\end{tabular}
\end{table}

\section{Problem Definition and Core Concepts}

\subsection{Intuition and Motivation}
Single linear layers can only express hyperplanar boundaries. Multi-Layer Perceptrons construct deep hierarchical representations where early layers extract low-level spatial or temporal primitives, and deeper layers combine these primitives non-linearly to partition complex decision manifolds.

\subsection{Formal Mathematical Definition}
\begin{definition}[MLP Forward Layer Recurrence]
Given input $a^{(0)} = x \in \mathbb{R}^{d_0}$, for each layer $l \in \{1, \dots, L\}$:
\begin{align}
z^{(l)} &= W^{(l)} a^{(l-1)} + b^{(l)} \label{eq:mlp_z} \\
a^{(l)} &= \begin{cases}
\sigma_{\text{hid}}\left(z^{(l)}\right) & \text{if } 1 \le l < L \\
\sigma_{\text{out}}\left(z^{(L)}\right) & \text{if } l = L
\end{cases} \label{eq:mlp_a}
\end{align}
The overall network function is the composite map:
\begin{equation}
f_{\text{MLP}}(x) = \left( \sigma_{\text{out}} \circ \mathcal{T}^{(L)} \circ \sigma_{\text{hid}} \circ \mathcal{T}^{(L-1)} \circ \dots \circ \sigma_{\text{hid}} \circ \mathcal{T}^{(1)} \right)(x)
\end{equation}
where $\mathcal{T}^{(l)}(v) \triangleq W^{(l)} v + b^{(l)}$.
\end{definition}

\section{Algorithmic Specification}

The computational execution implemented in \texttt{NnAlgoMlpFeedforward} is shown below.

\begin{algorithm}[H]
\caption{Multilayer Perceptron Forward Propagation}
\begin{algorithmic}[1]
\Require Input vector $x \in \mathbb{R}^{d_0}$, layer weights $\{W^{(1)}, \dots, W^{(L)}\}$, biases $\{b^{(1)}, \dots, b^{(L)}\}$
\Require Hidden activation $\sigma_{\text{hid}}$, output activation $\sigma_{\text{out}}$
\Ensure Final output vector $\hat{y} = a^{(L)}$, list of pre-activations $\{z^{(l)}\}$ and post-activations $\{a^{(l)}\}$
\State $L \gets \text{length}(W)$
\State $a \gets x$
\State $\text{all\_z} \gets \text{empty List}, \quad \text{all\_a} \gets [x]$
\For{$l \gets 1 \textbf{ to } L$}
    \State $d_{\text{out}} \gets \text{rows}(W[l]), \quad d_{\text{in}} \gets \text{cols}(W[l])$
    \State $z \gets \text{new Array of size } d_{\text{out}}$
    \For{$i \gets 1 \textbf{ to } d_{\text{out}}$}
        \State $\text{sum} \gets 0.0$
        \For{$j \gets 1 \textbf{ to } d_{\text{in}}$}
            \State $\text{sum} \gets \text{sum} + W[l][i][j] \cdot a[j]$
        \EndFor
        \State $z[i] \gets \text{sum} + b[l][i]$
    \EndFor
    \State $\text{all\_z}.\text{append}(z)$
    \State $\text{act\_fn} \gets (\sigma_{\text{out}} \text{ if } l = L \text{ else } \sigma_{\text{hid}})$
    \State $a \gets \text{new Array of size } d_{\text{out}}$
    \For{$i \gets 1 \textbf{ to } d_{\text{out}}$}
        \State $a[i] \gets \text{act\_fn}(z[i])$
    \EndFor
    \State $\text{all\_a}.\text{append}(a)$
\EndFor
\State \Return $\{\text{output}: a, \text{pre\_activations}: \text{all\_z}, \text{activations}: \text{all\_a}\}$
\end{algorithmic}
\end{algorithm}

\section{Theoretical Guarantees and Universal Approximation}

\begin{theorem}[Universal Approximation Theorem (Cybenko 1989, Hornik 1991)]
Let $\sigma: \mathbb{R} \to \mathbb{R}$ be any continuous non-polynomial activation function. Let $K \subset \mathbb{R}^{d_0}$ be any compact subset. For any continuous target function $f \in \mathcal{C}(K)$ and any error tolerance $\epsilon > 0$, there exists a single-hidden-layer MLP with finite hidden dimension $d_1 < \infty$ such that:
\begin{equation}
\sup_{x \in K} \left| \sum_{i=1}^{d_1} v_i \sigma(w_i^T x + b_i) - f(x) \right| < \epsilon
\end{equation}
\end{theorem}

\subsection{Limitations}
\begin{itemize}
    \item \textbf{Non-Constructive Bound}: The theorem guarantees existence, but gives no bound on training sample complexity or optimization tractability via gradient descent.
\end{itemize}

\section{Complexity Analysis}

\subsection{Time Complexity}
\begin{itemize}
    \item \textbf{Matrix Vector Multiplications}: Layer $l$ requires $d_l \cdot d_{l-1}$ multiply-accumulate operations.
    \item \textbf{Total Time}: $\Theta\left( \sum_{l=1}^L d_l d_{l-1} \right)$ FLOPs.
\end{itemize}

\subsection{Space Complexity}
\begin{itemize}
    \item \textbf{Forward Activation Cache}: Stores all intermediate pre- and post-activation vectors $\Theta\left( \sum_{l=0}^L d_l \right)$ for backpropagation.
\end{itemize}

\section{Worked Numerical Example}

Consider a 2-layer MLP with $d_0 = 2, d_1 = 2, d_2 = 1$, $\sigma_{\text{hid}} = \text{ReLU}$, $\sigma_{\text{out}} = \text{Identity}$:
\begin{equation}
x = \begin{bmatrix} 1.0 \\ -1.0 \end{bmatrix}, \quad W^{(1)} = \begin{bmatrix} 1.0 & 2.0 \\ -1.0 & 1.0 \end{bmatrix}, \quad b^{(1)} = \begin{bmatrix} 0.5 \\ 0.0 \end{bmatrix}
\end{equation}
\begin{equation}
W^{(2)} = \begin{bmatrix} 2.0 & -1.0 \end{bmatrix}, \quad b^{(2)} = [0.1]
\end{equation}

\begin{enumerate}
    \item \textbf{Layer 1 Pre-activations}:
    \begin{align}
    z_1^{(1)} &= 1.0(1.0) + 2.0(-1.0) + 0.5 = 1.0 - 2.0 + 0.5 = -0.50 \\
    z_2^{(1)} &= -1.0(1.0) + 1.0(-1.0) + 0.0 = -1.0 - 1.0 = -2.00
    \end{align}
    \item \textbf{Layer 1 Post-activations} (ReLU):
    \begin{align}
    a_1^{(1)} &= \max(0, -0.5) = 0.00 \\
    a_2^{(1)} &= \max(0, -2.0) = 0.00
    \end{align}
    \item \textbf{Layer 2 Output}:
    \begin{equation}
    z^{(2)} = 2.0(0.0) - 1.0(0.0) + 0.1 = 0.10 \implies a^{(2)} = 0.10
    \end{equation}
\end{enumerate}

\section{Numerical Considerations and Edge Cases}

\begin{itemize}
    \item \textbf{Dying ReLUs}: If large negative biases push all hidden unit pre-activations into $z < 0$, ReLU outputs identically zero and gradient flow dies.
    \item \textbf{Vanishing Gradients}: Sigmoid/tanh activations saturate at $|z| > 5$, decaying backward gradients exponentially with depth.
\end{itemize}

\begin{thebibliography}{9}
\bibitem{cybenko1989} G.~Cybenko, ``Approximation by superpositions of a sigmoidal function,'' \emph{Mathematics of Control, Signals, and Systems}, vol.~2, no.~4, pp.~303--314, 1989.
\bibitem{hornik1991} K.~Hornik, ``Approximation capabilities of multilayer feedforward networks,'' \emph{Neural Networks}, vol.~4, no.~2, pp.~251--257, 1991.
\end{thebibliography}

\end{document}
